\section{Tomamos colectivamente una decisión equivocada: complejidad, producto y personas}

El título del segundo capítulo es «¡Oh, tomamos colectivamente una decisión equivocada!». No se trata de un pequeño error aislado, sino de una muestra del ensayo y error constante de un equipo emprendedor entre la tecnología, el producto y la organización. Peak habla varias veces de la complejidad: la tecnología es compleja, el producto es complejo y las personas lo son todavía más. Comparada con la complejidad de un programa, la complejidad de las personas y de la organización es más difícil de controlar.

\subsection{Por qué las personas son más complejas que la IA}

Peak dice que no le gusta dirigir personas, porque la complejidad humana crece de forma casi exponencial a medida que aumenta el número de personas. Si un programa tiene patrones de diseño suficientemente buenos, su complejidad sigue siendo controlable; las personas y la organización, en cambio, generan problemas de comunicación, motivación, expectativas, emociones y distribución de autoridad y responsabilidad. Esto es especialmente importante para una empresa de agentes: la IA puede elevar la productividad individual, pero no resuelve por sí sola la complejidad organizacional.

\begin{warningbox}{Las herramientas de IA no eliminan por sí solas los problemas de la organización}
Aunque cada empleado use herramientas de IA, la organización sigue necesitando objetivos, colaboración, toma de decisiones, responsabilidad y flujo de información. La IA puede amplificar las capacidades individuales, pero también puede amplificar el desorden.
\end{warningbox}

\subsection{Cambios organizacionales en la era de la IA}

La sección anterior trató la complejidad de las personas y de la organización; esta sección examina cómo entra la IA en la organización. Peak considera que, en la era de la IA, los cambios organizacionales no son tan grandes como se suele imaginar, pero dentro de la organización habrá mucha más IA. Manus anima a sus empleados a usar todo tipo de productos de IA e incluso procura reembolsarles las herramientas de terceros, porque una empresa de IA debe, ante todo, lograr que sus empleados comprendan la vanguardia de la industria. Es un punto de vista muy práctico: no se trata de rediseñar primero el organigrama, sino de que cada persona obtenga primero el doble o diez veces más productividad.

\begin{importantbox}{El orden del cambio organizacional}
Primero hay que lograr que cada persona comprenda y use la IA, y después observar cómo cambian los flujos de trabajo y la estructura organizacional. Diseñar de forma prematura una «forma organizacional para la era de la IA» puede ser más lento que hacer que el equipo use realmente la IA.
\end{importantbox}

\subsection{Resumen de la sección}

Los errores y el ensayo y error de Manus no provienen solo de la dirección del producto, sino también de las personas y la organización dentro de un sistema complejo. Una empresa de agentes no solo tiene que resolver las capacidades del modelo, sino también cómo colaboran las personas, cómo mantener la simplicidad y cómo evitar que la complejidad de las funciones y de la organización la absorba.
